\chapter[The winners, part 2: models]{The winners, part 2:\\ models that know where to look}
\label{ch:winmodels}

\begin{tldr}
\begin{itemize}
\item Agents also drew \textbf{boxes}: around the \emph{anatomy} each finding concerns, and around visible
      \emph{lesions}. A detector learned to draw the anatomy boxes on every series.
\item \textbf{Box pooling}: each finding reads the image features only inside its own box, so ``lateral meniscus'' looks
      at the lateral meniscus.
\item Training-only extras: a \textbf{lesion head} and a \textbf{report head} (predict the report's meaning from the
      image). Free at test time.
\item Bigger input: all native slices, 384--512\,px, 5 series. A validation ruler checked against the leaderboard.
      One model: 0.954 private in \ca 4 minutes.
\end{itemize}
\end{tldr}

\section{Two kinds of boxes}

\begin{figure}[H]
\centering
\begin{tikzpicture}[scale=0.5]
  \kneesagittal
  % anatomy (ROI) box: the whole ACL region
  \draw[us, line width=1.6pt] (4.4,4.3) rectangle (6.6,6.8);
  \node[font=\scriptsize\bfseries, text=us, anchor=south west, fill=white, inner sep=1pt] at (4.4,6.85) {ACL anatomy box};
  % anatomy box: patellofemoral
  \draw[us, line width=1.6pt, dashed] (2.1,6.0) rectangle (4.3,9.1);
  \node[font=\scriptsize\bfseries, text=us, anchor=east, fill=white, inner sep=1pt] at (2.05,8.6) {PF OA box};
  % lesion boxes: bone bruise on femur and tibia
  \fill[them, opacity=0.35] (6.3,5.4) ellipse (0.55 and 0.4);
  \draw[them, line width=1.4pt] (5.7,4.95) rectangle (6.9,5.85);
  \fill[them, opacity=0.35] (6.6,4.0) ellipse (0.5 and 0.3);
  \draw[them, line width=1.4pt] (6.05,3.65) rectangle (7.15,4.35);
  \node[font=\scriptsize\bfseries, text=them, anchor=west, fill=white, inner sep=1pt] at (7.4,3.5) {lesion boxes};
\end{tikzpicture}\hspace{6mm}
\begin{minipage}[b]{0.48\linewidth}\small
\textcolor{us}{\textbf{Anatomy box}} (``region of interest'', ROI): where a finding is \emph{assessed}, whether or not
it is present. The ACL box covers the whole ligament; the PF OA box the kneecap and the femur's groove. A healthy knee
has these boxes too.\par\medskip
\textcolor{them}{\textbf{Lesion box}}: drawn tightly around a visible problem (here a bone bruise on femur and tibia),
on every slice where it is visible. Only sick knees have them.\par\bigskip
\end{minipage}
\caption{The two kinds of spatial labels. Anatomy boxes tell the model \emph{where to look}; lesion boxes tell it
\emph{what a problem looks like}.}
\label{fig:boxes}
\end{figure}

\section{From a few agent boxes to boxes everywhere}

Agents could only annotate part of the data: tascj's agents drew 3D anatomy boxes for 751 studies (21 sub-sessions in
parallel, managed by one main session); sakaku's drew 2D boxes slice by slice on 39,415 images from 3,823 studies; the
lesion agents boxed 3,616 lesions on 1,800 studies (111,029 boxes). A \term{detector} then learned to predict the
anatomy boxes for every slice of every series, including the test series.

\begin{figure}[H]
\centering
\begin{tikzpicture}[x=1mm,y=1mm,
  st/.style={draw=grid, fill=white, rounded corners=1.5mm, minimum height=13mm, align=center, font=\scriptsize, text width=30mm},
  ar/.style={-{Stealth[length=1.8mm]}, line width=0.7pt, muted}]
  \node[st, draw=them] (a) at (0,0) {\faRobot\ agent boxes\\ on \textbf{some} studies};
  \node[st, draw=accent, fill=keyfill] (b) at (42,0) {train a \textbf{box detector}\\ (a small CNN)};
  \node[st] (c) at (84,0) {boxes on \textbf{every}\\ slice of every series};
  \node[st] (d) at (126,0) {used by the\\ classifier};
  \draw[ar] (a) -- (b); \draw[ar] (b) -- (c); \draw[ar] (c) -- (d);
  \node[font=\scriptsize, text=muted, align=center] at (63,-14) {tascj: ConvNeXt-T per slice + 3D fusion, overlap with held-out boxes 0.71 (IoU)\\
    sakaku: DINOv3 ConvNeXt-B + DETR decoder, 0.95 mAP50};
\end{tikzpicture}
\caption{Expanding sparse annotations with a detector. tascj's detector also runs inside the submission; sakaku used
boxes for training only.}
\end{figure}

\begin{watchout}
\textbf{Missing is not negative.} A slice without a box can mean ``structure not visible here'' or ``nobody
annotated it''. sakaku stored three states per slice and finding: \emph{positive region}, \emph{reliable negative}
(checked: structure absent), \emph{ignore}. Treating every empty slice as negative would teach the model the
detector's misses.
\end{watchout}

\section{Box pooling}

\begin{figure}[H]
\centering
\begin{tikzpicture}[x=5mm,y=5mm]
  % feature map 10x10
  \foreach \i in {0,...,9} \foreach \j in {0,...,9} {
    \pgfmathsetmacro{\v}{10+mod(\i*7+\j*13,30)}
    \fill[us!\v] (\i,\j) rectangle ++(1,1);
  }
  \draw[white, line width=0.4pt] (0,0) grid (10,10);
  \draw[them, line width=1.8pt] (5,2) rectangle (8,5);
  \node[font=\scriptsize\bfseries, text=them, anchor=south west, fill=white, inner sep=1pt] at (5,5.05) {medial-meniscus box};
  \draw[third!80!black, line width=1.8pt, dashed] (1,2) rectangle (4,5);
  \node[font=\scriptsize\bfseries, text=third!70!black, anchor=north west, fill=white, inner sep=1pt] at (1,1.95) {lateral-meniscus box};
  \node[font=\scriptsize, align=center] at (5,-0.9) {feature map of one slice (each cell = a small patch of the image)};
  % outputs
  \draw[-{Stealth}, them, line width=1pt] (8.2,3.5) -- (12,3.5);
  \node[draw=them, fill=whyfill, rounded corners=1mm, font=\scriptsize, align=center, anchor=west] at (12,3.5) {average inside\\ the box $\rightarrow$\\ \textbf{Medial meniscus}};
  \draw[-{Stealth}, third!80!black, line width=1pt] (2.5,5.3) -- (2.5,8.5) -- (12,8.5);
  \node[draw=third!80!black, fill=wordfill, rounded corners=1mm, font=\scriptsize, align=center, anchor=west] at (12,8.5) {average inside\\ the box $\rightarrow$\\ \textbf{Lateral meniscus}};
\end{tikzpicture}
\caption{Box pooling on a coronal slice of a right knee (lateral side on the left, as in Figure~\ref{fig:anatomy}). Each finding averages only the features inside its own anatomy box (weighted by
how much of each cell the box covers), then one small linear layer turns that into the finding's score. A learned
\emph{in-box attention} later lets the model weight the most suspicious cells within the box more.}
\label{fig:boxpool}
\end{figure}

\begin{why}
Our model pooled whole slices (Figure~\ref{fig:attnpool}): it had to discover by itself which pixels belong to the
lateral meniscus. Box pooling hands it that knowledge. It also explains why the canonical side helped them and not us:
with boxes, ``medial'' and ``lateral'' are explicit regions; without boxes, the side was just one more thing to learn.
\end{why}

\section{Learning from boxes and reports during training only}

\begin{figure}[H]
\centering
\begin{tikzpicture}[x=1mm,y=1mm,
  st/.style={draw=grid, fill=white, rounded corners=1.5mm, minimum height=10mm, align=center, font=\scriptsize, text width=26mm},
  ar/.style={-{Stealth[length=1.8mm]}, line width=0.7pt, muted}]
  \node[st, draw=accent, fill=keyfill, minimum height=24mm] (bb) at (0,0) {backbone\\ feature maps};
  \node[st, draw=accent, line width=0.9pt] (main) at (45,14) {\textbf{box pooling}\\ $\rightarrow$ 12 findings};
  \node[st, dashed] (les) at (45,0) {\textbf{lesion head}\\ where are lesions?};
  \node[st, dashed] (rep) at (45,-14) {\textbf{report head}\\ predict the report's\\ embedding};
  \draw[ar] (bb) -- (main.west); \draw[ar, dashed] (bb) -- (les.west); \draw[ar, dashed] (bb) -- (rep.west);
  \node[font=\scriptsize, anchor=west, align=left] at (62,14) {soft labels (BCE) \hfill \textbf{used at test time}};
  \node[font=\scriptsize, anchor=west, align=left, text=muted] at (62,0) {lesion boxes (BCE per cell), weight 1};
  \node[font=\scriptsize, anchor=west, align=left, text=muted] at (62,-14) {report text $\rightarrow$ Qwen3-Embedding-8B (4,096 numbers);\\
    loss $= 1 - \cos(\text{predicted}, \text{actual})$, weight 0.3};
  \node[font=\scriptsize, text=muted] at (45,-24) {dashed = training only, removed for the submission};
\end{tikzpicture}
\caption{Auxiliary targets (tascj's model). They shape what the backbone learns but cost nothing at test time.}
\label{fig:auxheads}
\end{figure}

\begin{keyidea}
\textbf{The report-embedding head uses the report without needing it at test time.} An \term{embedding} is a list of
numbers that captures a text's meaning. Asking the image model to predict the report's embedding forces it to learn
everything the report talks about, not just the 12 findings. This needed no hosted LLM: an open-weights embedding model
was enough. It is the one idea that was fully open to us under our own rule.
\end{keyidea}

\section{The input: more of everything}

\begin{center}\small
\begin{tabularx}{\linewidth}{l>{\raggedright\arraybackslash}X>{\raggedright\arraybackslash}X}
\toprule
 & \textcolor{us}{\textbf{us}} & \textcolor{them}{\textbf{winners}} \\
\midrule
series & 3 fluid-sensitive & 5 roles (+ sagittal and coronal without fat suppression); one model type used every series as its own sample \\
slices & 24, resampled & up to 64 \textbf{native} slices (or 32 sampled, 30 windows at test time) \\
pixels & 140\,mm at 256\,px (0.55\,mm) & 140\,mm at 384--512\,px; joint crop 112\,mm at 384\,px (0.29\,mm) \\
image unit & 2.5D triplets & single slices (tascj) or triplets (sakaku) \\
left/right & mirror augmentation & left knees mirrored to right; side from DICOM header, else a side classifier (99.3\,\%) \\
\bottomrule
\end{tabularx}
\end{center}

\section{What each piece was worth}

\begin{figure}[H]
\centering
\begin{tikzpicture}
\begin{axis}[ybar, width=\linewidth, height=6.4cm, bar width=9mm, ymin=0.925, ymax=0.946,
    symbolic x coords={a,b,c,d,e,f}, xtick=data,
    xticklabels={{plain per-slice\\ ConvNeXt-T}, {+ box\\ pooling}, {+ 112\,mm\\ joint crop}, {+ in-box\\ attention},
                 {+ calibrated\\ labels}, {+ lesion and\\ report heads}},
    x tick label style={font=\scriptsize, align=center}, ylabel={gold macro AUC}, ylabel style={font=\footnotesize},
    ymajorgrids, nodes near coords, point meta=y, every node near coord/.append style={font=\scriptsize,
    /pgf/number format/fixed, /pgf/number format/precision=4, /pgf/number format/fixed zerofill},
    yticklabel style={/pgf/number format/fixed, /pgf/number format/precision=3}]
  \addplot[barthem] coordinates {(a,0.9319) (b,0.9330) (c,0.9386) (d,0.9396) (e,0.9410) (f,0.9432)};
\end{axis}
\end{tikzpicture}
\caption{tascj's step-by-step test at 256\,px, one component at a time, on the 58 gold studies (5 fold models). Note
the scale: the whole architecture adds \textbf{+0.011}. The starting point, 0.932, is already far above our 0.904,
because of the labels and the input.}
\label{fig:winablation}
\end{figure}

\section{Training}

\begin{center}\small
\begin{tabularx}{\linewidth}{l>{\raggedright\arraybackslash}X>{\raggedright\arraybackslash}X}
\toprule
 & tascj & sakaku \\
\midrule
backbones & ConvNeXt-T, MambaOut-B+, MaxxViT-V2-B & MaxViT-V2-B, CoAtNet-rMLP-2 \\
learning rate & $5\cdot10^{-5}$, constant after warm-up & backbone $8.9\cdot10^{-5}$ (layer-wise decay 0.9), new layers $10^{-3}$, one-cycle \\
epochs & 8, EMA 0.999 & up to 30 for tests; final 19, average of epochs 15--19 \\
loss & soft BCE + lesion + report & soft BCE + 0.05 ranking loss + box losses \\
precision & fp16 + \code{torch.compile} & bf16 + activation checkpointing \\
\bottomrule
\end{tabularx}
\end{center}

\begin{wordcheck}
\term{EMA} (exponential moving average): keep a slowly updated average copy of the weights and use that copy;
it is usually a little better than the latest weights. \term{Layer-wise decay}: early backbone layers learn more
slowly than late ones.
\end{wordcheck}

\begin{innumbers}
sakaku found that raising the backbone learning rate about 10$\times$ improved the public score from 0.936 to 0.944,
while the 58-study gold score went \emph{down} (0.928 $\rightarrow$ 0.922). That disagreement is what pushed the team to
build a better ruler.
\end{innumbers}

\section{A ruler they could trust}

Instead of the noisy 58 gold studies, they scored models on out-of-fold predictions for thousands of report-labelled
studies, against their best agent labels. To decide \emph{which} ruler to trust, they collected 20--30 leaderboard
scores and kept the one that tracked the leaderboard best.

\begin{figure}[H]
\centering
\begin{tikzpicture}[x=1mm,y=1mm]
  \node[font=\small] at (0,0) {$\displaystyle \text{SoftAUC} = \frac{\sum_{i \neq j} q_i\,(1-q_j)\cdot[\,\text{pred}_i > \text{pred}_j\,]}{\sum_{i\neq j} q_i\,(1-q_j)}$};
  \node[font=\scriptsize, align=left, anchor=west] at (48,0) {$q$ = reference soft label.\\
     A pair counts as much as ``$i$ is sick and $j$ is healthy''\\ is likely; ties get half credit.};
\end{tikzpicture}
\caption{\term{SoftAUC}: AUC against soft reference labels, without rounding them to 0/1 (sakaku). It keeps the
labels' uncertainty in the measurement.}
\end{figure}

\section{The final ensemble}

\begin{center}\small
\begin{tabularx}{\linewidth}{>{\raggedright\arraybackslash}X l l c}
\toprule
Model (all 384\,px) & Labels & Copies & Weight \\
\midrule
\multicolumn{4}{l}{\textit{tascj's half}} \\
ConvNeXt-T (alone: 0.967 public / 0.958 private) & trained-guide Astra & 5 folds & 2 \\
MambaOut-B+ & calibrated Astra (+ extra heads) & full data $\times$ 3 seeds & 1 \\
MaxxViT-V2-B & same & full data $\times$ 3 seeds & 1 \\
ConvNeXt-T, 6-channel & mean of two Astra versions & full data $\times$ 3 seeds & 1 \\
MambaOut-B+, 6-channel & same & full data $\times$ 3 seeds & 1 \\
\multicolumn{4}{l}{\textit{sakaku's half (alone: 0.967 public / 0.957 private)}} \\
MaxViT-V2-B & Mix5 (five sources) & full data & 1/3 \\
MaxViT-V2-B & focus-guided Astra & full data & 1/3 \\
CoAtNet-rMLP-2 & Mix5 & full data & 1/3 \\
\midrule
\textbf{Team, 50/50} & & & \textbf{0.969 / 0.960} \\
\bottomrule
\end{tabularx}
\end{center}

\begin{wordcheck}
\term{6-channel input}: instead of 3 neighbouring slices of one series, each input image has 6 channels, one per
(plane $\times$ fat suppression) kind, all resampled in 3D onto the current slice. The model sees the same spot in
several sequences at once. According to tascj this added more than switching backbones did.
\end{wordcheck}

\section{The efficiency notebook: 0.954 in four minutes}

\begin{itemize}
\item One ConvNeXt-T at 256\,px, trained-guide labels, all training data, one checkpoint, plus the box detector.
\item Inference compiled with \term{TensorRT} (NVIDIA's optimiser), one process per T4, DICOM reading in C++, and only
      the fluid-sensitive series read (41\,\% less data).
\item With 1,200 training studies laid out as a test set: \textbf{230 seconds} for the whole notebook.
      \textbf{0.964 public / 0.954 private}.
\end{itemize}

\begin{innumbers}
With the efficiency formula (Chapter~\ref{ch:competition}), if the hidden-test run was as fast as their own test:
$-2.16 \times 0.954 + 0.07/9 \approx -2.05$. A 1-hour notebook
would need about 0.05 more AUC to match it. Kaggle's disk speed varied from \ca 30 to 500+\,MB/s with platform load, so
they submitted when the platform was quiet.
\end{innumbers}

\section{How they worked}

Both write-ups describe a workflow driven by AI coding agents: one agent was asked to develop and improve the
labelling instructions against the gold studies, and tascj notes that most model ideas came from asking Claude
(Opus 5.5) to propose ways to use the boxes and to run the experiments. The humans chose the direction, built the
measuring stick, and decided what to trust.

\begin{checkpoint}
\begin{itemize}
\done I can explain the difference between an anatomy box and a lesion box.
\done I can explain box pooling and why it helps side-specific findings.
\done I know what auxiliary heads are and why they are free at test time.
\quiz Q13.1: Why is predicting the report's embedding useful when the test set has no reports?
\quiz Q13.2: The team's whole architecture added +0.011 on gold. Where did the rest of their lead over us come from?
\quiz Q13.3: Why did they choose their validation ruler by correlation with the leaderboard instead of using gold?
\end{itemize}
\end{checkpoint}
